\begin{table}[!t]
\centering
\caption{Ways in which a reply can act on its host, and our defense against each one.}
\label{tab:threats}
\footnotesize
\renewcommand{\arraystretch}{1.12}
\begin{tabular}{@{}>{\raggedright\arraybackslash}p{3.85cm}>{\raggedright\arraybackslash}p{3.95cm}@{}}
\toprule
Threat & Defense \\
\midrule
A file name with a directory part, such as \code{../../.bashrc} & Refused before any write. \\
A check that reads the home folder or the API key & Neither exists in the sandbox. \\
A check that sends data over the network & Only a loopback interface exists. \\
A check that changes host files or leaves a link for a later write & A read-only system, and a working folder that disappears with each check. \\
A check that runs forever or uses too much memory & Limits on time, memory, file size and open files (\cref{tab:limits}). \\
A check that fills memory with many files & Open in the harness. A cgroup stops it (\cref{tab:cgroup}). \\
A check that starts processes in a loop & Open in the harness. A cgroup stops it (\cref{tab:cgroup}). \\
A check that uses a kernel bug & Open. We load no seccomp filter. \\
Output that carries instructions to the model & Open. We return its output unchanged. \\
A check that tests nothing & Open. A model writes its own check. \\
\bottomrule
\end{tabular}
\end{table}
